The Open-Source Software ecosystem relies heavily on package managers such as Node Package Manager (npm), 
which hosts over two million libraries. 
However, its open nature makes it vulnerable to Software Supply Chain Attacks (SSCAs), 
where malicious code is injected into legitimate components. 
Existing detection tools share two limitations: 
they analyze only the most recent release, lacking a longitudinal perspective, 
and incur high operational costs through sandbox execution or commercial Large Language Model inference.

This thesis proposes a metric-based approach to identify anomalous npm releases 
by tracking their evolution across versions. 
We first study four SSCAs that affected the npm ecosystem in 2025: 
a malicious Windows library injection, a cryptocurrency-stealing attack, 
and two \emph{Shai-Hulud} self-propagating worm variants. 
From these, we define five metrics: 
\emph{File Types}, \emph{Package Size}, \emph{Longest Line Length}, \emph{Unique Hexadecimal Values}, and \emph{Unique Ethereum Addresses}.
We evaluate these metrics against a goodware baseline of 629 popular packages across their 20 most recent versions, 
and a malware collection of 20 compromised releases. 

Results show that each attack is detectable by at least one metric, 
but no single one identifies all four simultaneously, 
highlighting the value of combining lightweight measures 
as a practical complement to semantically rich approaches.